\section{Анализ предметной области}
 
\subsection{Понятие и принципы онлайн-аренды вещей}
 
Электронная коммерция представляет собой совокупность коммерческих и финансовых операций, осуществляемых посредством глобальных компьютерных сетей, а также комплекс сопутствующих бизнес-процессов, обеспечивающих их проведение. В рамках данной сферы онлайн-аренда материальных объектов сформировалась как самостоятельный сегмент, характеризующийся временной передачей права пользования вещью через интернет-платформу на возмездной основе без перехода права собственности к пользователю.
 
Концептуальной основой онлайн-аренды служит модель совместного потребления (sharing economy), предполагающая эффективное распределение недозагруженных материальных ресурсов посредством цифровых посредников. С инженерной точки зрения реализация данной концепции требует создания многоуровневой программно-информационной системы, центральным звеном которой выступает серверная часть, несущая ответственность за обработку бизнес-логики, персистентное хранение данных, управление транзакциями и обеспечение информационной безопасности.
 
Принята следующая классификация субъектов электронной коммерции по типу участников:
\begin{itemize}
    \item C2C (Consumer-to-Consumer) — взаимодействие между конечными потребителями;
    \item B2C (Business-to-Consumer) — взаимодействие коммерческой организации с конечным потребителем;
    \item B2B (Business-to-Business) — взаимодействие хозяйствующих субъектов между собой;
    \item B2G (Business-to-Government) — взаимодействие коммерческих структур с органами государственного управления;
    \item P2P (Peer-to-Peer) — децентрализованное взаимодействие равноправных участников без выделенного посредника.
\end{itemize}
 
В сегменте аренды вещей доминируют модели C2C и B2C, что непосредственно определяет требования к ролевой модели и разграничению прав доступа в проектируемой серверной части.
 
Онлайн-бронирование представляет собой процесс резервирования товара или услуги в интерактивном режиме реального времени посредством сети Интернет. Данный механизм широко применяется при резервировании номерного фонда гостиниц, приобретении транспортных билетов, бронировании столиков в заведениях общественного питания и учреждениях культуры. Применительно к аренде вещей онлайн-бронирование приобретает дополнительную специфику: серверная часть системы обязана обеспечивать управление временной доступностью каждого объекта, предотвращение конкурентных резервирований, автоматизированный расчёт стоимости, а также управление финансовыми обязательствами сторон в течение всего жизненного цикла сделки.
 
Использование онлайн-платформ аренды обеспечивает следующие преимущества для участников, каждое из которых транслируется в конкретные требования к серверным компонентам системы:
 
\begin{enumerate}
    \item Экономическая эффективность. Арендаторы получают доступ к дорогостоящим или редко используемым объектам без необходимости их приобретения в собственность; арендодатели извлекают доход от временно незадействованного имущества. Серверная часть обеспечивает автоматизированный расчёт итоговой стоимости с учётом комиссионного вознаграждения платформы.
    \item Доступность и масштабируемость. Поиск предложений, оформление бронирования и проведение оплаты доступны в непрерывном режиме из любой географической точки при наличии подключения к сети Интернет. Архитектура серверной части должна обеспечивать стабильную обработку запросов при значительных пиковых нагрузках.
    \item Прозрачность и верифицируемость сделок. Встроенная система рейтингов и отзывов формирует репутационный профиль каждого участника, снижая информационную асимметрию при заключении сделок. Серверная часть обеспечивает хранение, агрегирование рейтинговых данных и их учёт в алгоритмах ранжирования объявлений.
\end{enumerate}
 
Наряду с преимуществами функционирование платформы сопряжено с рисками, которые подлежат учёту на стадии проектирования серверной части:
 
\begin{enumerate}
    \item Несоответствие объекта аренды заявленным характеристикам. Ввиду отсутствия возможности физической инспекции предмета до завершения бронирования существует вероятность расхождения между реальным состоянием объекта и его описанием в карточке объявления. Для снижения данного риска серверная часть должна реализовывать валидацию и модерацию загружаемых медиаматериалов.
    \item Риски, связанные с передачей и возвратом имущества. Организация логистики влечёт дополнительные издержки, а в процессе доставки или эксплуатации возможны утрата либо повреждение предмета аренды. Серверная часть фиксирует статусы передачи и приёмки, формируя доказательную базу для урегулирования споров.
    \item Юридические и финансовые разногласия. При отсутствии встроенных механизмов обеспечения сделки урегулирование споров о возмещении ущерба или возврате депозита существенно затрудняется. Реализация на стороне сервера механизма условного удержания денежных средств (escrow) и автоматической генерации договора аренды снижает данные риски.
    \item Угрозы информационной безопасности. Серверная часть является основным объектом атак, направленных на хищение персональных данных пользователей и осуществление несанкционированных финансовых операций. Обеспечение безопасности достигается посредством применения JWT-аутентификации, шифрования канала передачи данных по протоколу HTTPS, а также использования сертифицированных платёжных шлюзов.
\end{enumerate}
 
\subsection{Классификация онлайн-площадок для аренды}
 
Современный рынок сервисов онлайн-аренды характеризуется значительной неоднородностью предлагаемых решений. Систематизация существующих платформ по ключевым признакам позволяет обоснованно определить место проектируемой системы в данном сегменте и сформулировать архитектурные требования к её серверной части.
 
\subsubsection{По составу предлагаемых товарных категорий}
 
По степени предметной специализации выделяются два типа платформ: универсальные и нишевые (специализированные).
 
Универсальные платформы агрегируют объявления по широкому спектру категорий: электроника, строительный инструмент, спортивный инвентарь, предметы одежды и т.д. Конкурентным преимуществом таких систем служат масштаб аудитории и объём предложений. Вместе с тем в части серверного проектирования универсальность обусловливает необходимость поддержки полиморфной структуры атрибутов объявлений, что усложняет схему базы данных и логику серверной валидации.
 
Нишевые платформы сосредоточены на одной или нескольких смежных товарных категориях. Данное ограничение позволяет применить фиксированную реляционную схему для хранения атрибутов объявлений, упростить серверную валидацию и обеспечить более высокую релевантность поисковой выдачи. Проектируемая система тяготеет к универсальной модели, что определяет соответствующие требования к архитектуре модуля каталога.
 
\subsubsection{По модели участия платформы в транзакции}
 
В зависимости от степени вовлечённости платформы в процесс заключения и сопровождения сделки различают три функциональные модели.
 
Модель классифайда предполагает, что платформа выполняет исключительно функцию доски объявлений: публикует предложения арендодателей, тогда как условия оплаты и передачи имущества участники согласовывают самостоятельно за пределами системы. Серверная часть в данном случае ограничивается функциями хранения объявлений и базовой аутентификации пользователей.
 
Модель полного сопровождения сделки предусматривает, что платформа принимает на себя обработку и временное удержание платежей, генерацию юридически значимых документов, информирование сторон о смене статусов и участие в урегулировании споров, взимая комиссионное вознаграждение с каждой транзакции. Именно данная модель положена в основу проектируемой системы. Её реализация предъявляет наиболее высокие требования к серверной архитектуре: необходимы самостоятельные функциональные модули для обработки платежей, генерации договоров, WebSocket-коммуникации и управления уведомлениями.
 
Подписочная модель предполагает внесение пользователем периодического платежа за право неограниченного доступа к аренде вещей из заранее сформированного каталога. Данная модель наиболее востребована в категориях детских товаров и одежды.
 
\subsubsection{По типу арендодателей}
 
По составу участников сделки на стороне предложения выделяются две основные модели. Модель C2C характеризуется тем, что арендодателями выступают физические лица, предоставляющие во временное пользование объекты личной собственности. Модель B2C предполагает участие коммерческих организаций, осуществляющих профессиональную деятельность по прокату имущества и использующих платформу в качестве дополнительного канала привлечения клиентов. Поскольку проектируемая система ориентирована на поддержку обеих моделей, серверная часть должна реализовывать гибкое ролевое разграничение с возможностью совмещения ролей арендодателя и арендатора в рамках одной учётной записи.
 
\subsection{Основные бизнес-процессы системы аренды}
 
\subsubsection{Регистрация и верификация пользователей}
 
Обеспечение взаимного доверия участников платформы и снижение рисков мошенничества и нецелевого использования имущества достигается посредством многоуровневой процедуры регистрации и верификации. Серверная часть реализует последовательность этапов: создание учётной записи с подтверждением адреса электронной почты через одноразовый верификационный код, опциональную привязку платёжного инструмента, а в установленных случаях — верификацию личности посредством интеграции с государственными информационными сервисами. Верификационные данные сохраняются в базе данных в зашифрованном виде с жёстким разграничением прав доступа. Архитектура ролевой модели допускает совмещение функций арендодателя и арендатора в рамках единой учётной записи.
 
\subsubsection{Публикация и поиск предложений}
 
Для размещения объекта аренды арендодатель формирует карточку объявления, включающую наименование, категорию, техническое описание, фотографические материалы, ценовые условия и периоды доступности. Серверная часть обеспечивает приём и хранение медиафайлов, индексирование текстовых атрибутов для реализации полнотекстового поиска, а также фильтрацию результатов по заданным критериям: категории, ценовому диапазону и датам доступности.
 
\subsubsection{Бронирование и совершение сделки}
 
Процесс оформления аренды автоматизируется на основе встроенного серверного механизма управления бронированиями, исключающего возможность одновременного резервирования одного объекта несколькими пользователями. После выбора арендатором периода аренды серверная часть автоматически вычисляет итоговую стоимость с учётом комиссионного вознаграждения платформы. Проведение расчётов осуществляется исключительно через платформу: денежные средства временно блокируются на стороне сервера до получения подтверждений от обеих сторон о выполнении своих обязательств. Данный механизм гарантирует арендодателю получение вознаграждения, а арендатору — передачу предмета аренды в состоянии, соответствующем условиям договора.
 
\subsubsection{Система рейтингов и отзывов}
 
Механизм репутационной оценки участников является одним из ключевых инструментов снижения информационной асимметрии в C2C-сервисах. По завершении каждой сделки стороны имеют возможность выставить оценку и оставить текстовый отзыв. Серверная часть сохраняет поступившие данные, пересчитывает агрегированный рейтинг пользователя и учитывает его в алгоритме ранжирования объявлений в результатах поиска. Накопленная репутация оказывает непосредственное влияние на видимость предложений и готовность других участников платформы вступать в сделки с данным пользователем.
 
\subsection{Компоненты серверной части программно-информационной системы}
 
Серверная часть онлайн-площадки представляет собой совокупность специализированных программных модулей, реализованных по принципу разделения ответственности и предоставляющих свои функции клиентскому приложению через унифицированный REST API. Взаимодействие модулей осуществляется посредством внутренних интерфейсов; персистентное хранение данных обеспечивается реляционной СУБД.
 
\subsubsection{Модуль управления пользователями}
 
Модуль реализует функции регистрации учётных записей, аутентификации на основе JSON Web Token (JWT) с поддержкой механизма обновления токенов по истечении срока их действия, а также авторизации на основе атрибутивной ролевой модели доступа. В базе данных персистируются профили пользователей, история транзакций, рейтинговые показатели, текстовые отзывы и документы, полученные в ходе верификации личности.
 
\subsubsection{Модуль каталога и поиска}
 
Модуль экспонирует набор API-endpoints для выполнения операций CRUD над записями объявлений. Серверная логика включает загрузку, валидацию и хранение медиафайлов посредством выделенного файлового хранилища, построение и обновление поискового индекса, а также обработку поисковых запросов с применением комбинированной фильтрации по категориям, ценовому диапазону и временным периодам доступности.
 
\subsubsection{Модуль бронирования и управления доступностью}
 
Функциональная ответственность модуля состоит в управлении временной доступностью объектов аренды и обеспечении атомарности операций резервирования с целью исключения конкурентных конфликтов. Серверная часть поддерживает актуальный календарь занятости в разрезе каждого объявления, отслеживает конечный автомат состояний заявки на бронирование (создана, подтверждена, оплачена, завершена, отменена) и инициирует отправку уведомлений при каждом переходе между состояниями.

 
\subsubsection{Модуль аналитики и отчётности}
 
Модуль агрегирует данные о завершённых транзакциях и предоставляет арендодателям персональный дашборд с показателями доходов и статистикой просмотров объявлений. Для арендаторов формируется история бронирований с детализацией по статусам. На уровне администрирования платформы реализовано формирование сводных отчётов о совокупном обороте и количестве совершённых сделок за произвольно заданный временной период.
 
\subsection{Правовое регулирование онлайн-платформ аренды в Российской Федерации}
 
Проектирование серверной части онлайн-площадки осуществляется с учётом действующей нормативно-правовой базы, определяющей правовой статус участников арендных отношений, требования к обработке персональных данных и порядку проведения финансовых операций.
 
Гражданский кодекс Российской Федерации, глава 34 «Аренда», закрепляет существенные условия договора аренды, права и обязанности арендодателя и арендатора, а также ответственность за недостатки переданного в пользование имущества\cite{zakon2}. Совокупность данных норм определяет обязательный состав реквизитов и структуру разделов документов, автоматически генерируемых серверным модулем.
 
Закон РФ от 07.02.1992 № 2300-1 «О защите прав потребителей» распространяется на правоотношения с арендаторами, использующими предмет аренды в личных или бытовых целях. В соответствии с данным законом арендатор наделён правом на получение исчерпывающей и достоверной информации об услуге, на расторжение договора в любой момент с возмещением арендодателю фактически понесённых расходов, а также на компенсацию убытков, возникших вследствие недостатков оказанной услуги. Серверная часть должна обеспечивать неизменяемое хранение условий сделки в том виде, в котором они были приняты сторонами\cite{zakon1}.
 
Федеральный закон от 27.07.2006 № 152-ФЗ «О персональных данных» обязывает платформу как оператора персональных данных обеспечить законность их сбора и обработки, а также надёжную техническую защиту от несанкционированного доступа. Персональные данные пользователей — паспортные сведения, контактная информация, платёжные реквизиты — подлежат хранению в зашифрованном виде с разграничением доступа на уровне серверной части\cite{zakon4}. Интеграция с платёжными сервисами, прошедшими сертификацию по стандартам PCI~DSS или ISO-20022, обеспечивает соответствие требованиям к безопасности финансовых транзакций.
 
Федеральный закон от 27.11.2018 № 422-ФЗ устанавливает специальный налоговый режим «Налог на профессиональный доход», применение которого доступно арендодателям --- физическим лицам, не имеющим статуса индивидуального предпринимателя, в качестве правового основания для легального получения дохода от сдачи имущества в аренду\cite{zakon3}.
 
\subsection{Анализ существующих решений и аналогов}
 
С целью выработки обоснованных технических требований к проектируемому серверному решению проведён сравнительный анализ трёх наиболее востребованных отечественных платформ в сегменте онлайн-аренды: Rentomania, Авито и Арендую.ру. Оценка производилась по критериям, непосредственно связанным с функциональностью серверной части. Результаты сопоставления представлены в таблице~\ref{concurent:table}.
 
\begin{xltabular}{\textwidth}{|X|X|X|X|X|}
	\caption{Сравнение существующих решений и аналогов\label{concurent:table}}\\ \hline
	Параметр & \centrow Rentomania & \centrow Авито & \centrow Арендую.ру & \centrow Проекти-\ руемая онлайн площадка\\ \hline
	\endfirsthead
	\continuecaption{Продолжение таблицы \ref{concurent:table}}
	~ Параметр & \centrow Rentomania & \centrow Авито & \centrow Арендую.ру & \centrow Проекти-\ руемая онлайн площадка\\ \hline
	\finishhead
	\leftrow Размещение объявлений арендодателем & \centrow + & \centrow + & \centrow + & \centrow + \\ \hline 
	\leftrow Поиск и фильтрация товаров & \centrow + & \centrow + & \centrow + & \centrow + \\ \hline
	\leftrow Онлайн-бронирование & \centrow + & \centrow -- & \centrow -- & \centrow + \\ \hline
	\leftrow Встроенная онлайн-оплата & \centrow + & \centrow + & \centrow -- & \centrow + \\ \hline
	\leftrow Система рейтинга и отзывов & \centrow + & \centrow + & \centrow -- & \centrow + \\ \hline
	\leftrow Чат между арендатором и арендодателем & \centrow + & \centrow + & \centrow -- & \centrow + \\ \hline
	\leftrow Адаптивная вёрстка для мобильных устройств & \centrow + & \centrow + & \centrow -- & \centrow + \\ \hline
	\leftrow Верификация пользователей & \centrow + & \centrow -- & \centrow -- & \centrow + \\ \hline
	\leftrow Панель аналитики для арендодателя & \centrow -- & \centrow -- & \centrow -- & \centrow + \\ \hline
	\leftrow Понятный интерфейс & \centrow -- & \centrow + & \centrow -- & \centrow + \\ \hline
	\leftrow Уведомления о статусе аренды & \centrow + & \centrow -- & \centrow -- & \centrow + \\ \hline
\end{xltabular}